%%
%% Test: Cover System — v1.1
%% Stage: cover refactor (stages 1-4)
%% 
%% Purpose: Verify cover system:
%%   - Front cover (\makecover)
%%   - Back cover (\makebackcover)
%%   - Color aliases (cover* → matin*)
%%   - Repository variable (\matin@repository)
%%   - Version variable (\matinbookversion)
%%   - Requires 2 compilation passes
%% 

\documentclass{matinbook}

\title{تست سیستم جلد — نسخه ۱.۱}
\author{تیم MatinBook}
\date{۱۴۰۵}

\booktitle{تست جلد}

\begin{document}

%----------------------------------------
% Front Cover
%----------------------------------------
\makecover
    {عنوان کتاب نمونه}
    {زیرعنوان کتاب نمونه}
    {نام نویسنده}
    {نام ناشر}

%----------------------------------------
% Main Content
%----------------------------------------
\chapter{بررسی سیستم جلد}
\label{chap:test-cover}

این سند، سیستم جلد \texttt{mb-theme-cover} را بررسی می‌کند.

%----------------------------------------
\section{رنگ‌های جلد}
\label{sec:cover-colors}

رنگ‌های جلد از \texttt{matin*} مشتق شده‌اند:

\begin{itemize}
    \item \textcolor{coverbg}{\textbf{coverbg}} → \texttt{matincream}
    \item \textcolor{coverprimary}{\textbf{coverprimary}} → \texttt{matindarkblue}
    \item \textcolor{coveraccent}{\textbf{coveraccent}} → \texttt{matingreen}
    \item \textcolor{coveraccent2}{\textbf{coveraccent2}} → \texttt{matinorange}
    \item \textcolor{coveraccent3}{\textbf{coveraccent3}} → \texttt{matinblue}
    \item \textcolor{covergray}{\textbf{covergray}} → \texttt{matinmediumgray}
    \item \textcolor{coverlight}{\textbf{coverlight}} → \texttt{matinlightgray}
    \item \textcolor{coverdark}{\textbf{coverdark}} → \texttt{matindarkgray}
\end{itemize}

%----------------------------------------
\section{متغیرهای جلد}
\label{sec:cover-variables}

\begin{itemize}
    \item \texttt{\textbackslash matinbookversion} = \matinbookversion
    \item \texttt{\textbackslash repository} = \repository
    \item \texttt{\textbackslash covertag} = \covertag
    \item \texttt{\textbackslash coverabout} = \coverabout
    \item \texttt{\textbackslash coversource} = \coversource
\end{itemize}

%----------------------------------------
\section{تست برچسب‌های TikZ}
\label{sec:tikz-labels}

\begin{latin}
\begin{tikzpicture}
    \fill[coverprimary] (0,0) rectangle (2,1);
    \fill[coveraccent] (2.5,0) rectangle (4.5,1);
    \fill[coveraccent2] (5,0) rectangle (7,1);
    \node[color=white] at (1,0.5) {A};
    \node[color=white] at (3.5,0.5) {B};
    \node[color=white] at (6,0.5) {C};
\end{tikzpicture}
\end{latin}

%----------------------------------------
\section{نتیجه‌گیری}
\label{sec:conclusion}

اگر این سند بدون خطا کامپایل شود:

\begin{itemize}
    \item ✅ \texttt{\textbackslash makecover} کار می‌کند
    \item ✅ \texttt{\textbackslash makebackcover} کار می‌کند
    \item ✅ رنگ‌های \texttt{cover*} به \texttt{matin*} نگاشت شده‌اند
    \item ✅ \texttt{\textbackslash matin@repository} کار می‌کند
    \item ✅ \texttt{\textbackslash matinbookversion} کار می‌کند
    \item ✅ \texttt{\textbackslash covergeometric} با TikZ کار می‌کند
    \item ✅ جلد نیاز به ۲ بار کامپایل دارد
\end{itemize}

\textbf{وضعیت تست:} قبول

%----------------------------------------
% Back Cover
%----------------------------------------
\makebackcover{%
    این یک کتاب نمونه است که با کلاس \texttt{MatinBook} حروف‌چینی شده است.
    هدف این کتاب، آموزش برنامه‌نویسی و ریاضیات به زبان فارسی است.
}

\end{document}
